\documentclass[10pt,a4paper]{amsart}
\usepackage[margin=19mm]{geometry}
\usepackage{amsmath,amssymb,mathtools}
\usepackage[scr=rsfs,cal=euler]{mathalfa}
\usepackage{xfrac}
\usepackage[unicode,hidelinks,bookmarks=false]{hyperref}
\newcommand{\ie}{i.e.\ }
\newcommand{\cf}{cf.\ }
\newcommand{\resp}{respectively\ }
\newcommand{\Subsection}{\subsection}
\providecommand{\nobreakdash}{\nobreak\hbox{-}}
\setlength{\emergencystretch}{2em}
\multlinegap0pt
\usepackage{amssymb}
\usepackage{mathtools}
\usepackage{csquotes}
\usepackage{xcolor}
\newtheorem{lem}{Lemma}
\newcommand{\ud}{\mathrm{d}}
\title{\textbf{On the regularity of solutions to the Hamilton-Jacobi equations for the $N$-body problem}}
\author{Diego Berti, Davide Polimeni, Susanna Terracini}
\date{Source: arXiv:2507.19170v2 (CC BY 4.0)}
\begin{document}
\maketitle

\noindent\emph{Source and licence.} \textbf{On the regularity of solutions to the Hamilton-Jacobi equations for the $N$-body problem}. Version arXiv:2507.19170v2 (CC BY 4.0), distributed by arXiv under the Creative Commons Attribution 4.0 International licence, http://creativecommons.org/licenses/by/4.0/, as recorded in the arXiv licence metadata for this exact version and shown on the abstract page https://arxiv.org/abs/2507.19170v2. Full licence notice: \texttt{licenses/arxiv-2507.19170-CC-BY-4.0.txt}.\par\smallskip

\noindent\textit{Editorial scope.} Counted unit: the article's lem lem:Eq\_u\_T (source0.tex lines 1477--1497). The article's complete original proof of it is delivered with it (source0.tex lines 1535--1628, one \texttt{proof} environment of 94 lines, which the article places away from the statement). No mathematics is written, repaired or re-derived: the statement and the proof are the article's own lines, in the article's own notation, and no retained line is substituted or re-flowed.  Retained uncounted as the source-local context the proof needs: lines 1289--1289; lines 1409--1415; lines 1440--1447; lines 1454--1456.  Nothing was omitted from any retained span, so the package carries no omission record.  No retained line contains a citation, so the wrapper carries no bibliography and no bibliography entry.  No retained line contains a figure environment, an image-inclusion command or drawing code, so no image is delivered and the package declares no asset dependency.  Editorial formatting only, in addition: an AMS \texttt{amsart} preamble that restates the article's own declarations and macros used by the retained lines, namely the environments \texttt{lem} on separate counters, together with the standard packages those lines need.  The retained environments print in this document as Lemma 1 in order of appearance, whereas the article prints the same items with the numbering of its own full document.  The complete native source of the article as distributed in the arXiv e-print for this exact version is bundled unchanged as \texttt{sources/182193/source0.tex}.
\subsection{Semiconcavity with linear modulus of the value function}\label{sez_semiconcavity}

\begin{equation}\label{eq:u(T,x)_vers1}
\begin{split}
u(T,x)&= \min_{\eta \in H^1([1,T]):\,\eta(1)=x}\int_1^T L(\eta(t),\dot \eta(t))\,\ud t - \langle \dot r_0(T), \eta(T)\rangle_\mathcal{M}
\\
&=\int_1^T L(\gamma(t),\dot{\gamma}(t))\,\ud t - \langle \dot r_0(T), \gamma(T)\rangle_\mathcal{M},
\end{split}
\end{equation}

\begin{equation}\label{eq:u(T,x)_vers2}
\begin{split}
u(T,x) &=\int_1^T L(\hat\gamma(t),\dot{\hat\gamma}(t))\, \ud t -\langle \dot r_0(T), \hat\gamma(1)\rangle_\mathcal{M}
\\
&= \min_{\eta \in H^1([1,T]):\, \eta(T)=x}\int_1^T L(\eta(t),\dot\eta(t))\, \ud t -
\langle \dot r_0(T), \eta(1)\rangle_\mathcal{M}.
\end{split}
\end{equation}

    \begin{equation}\label{eq:dpp}
        u(T,x) \leq u(T',\overline{\gamma}(T')) + \int_{T'}^{T} L(\overline{\gamma}(t),\dot{\overline{\gamma}}(t))\ \ud t - \langle\dot{r}_0(T)-\dot{r}_0(T'),\Tilde{\gamma}(1)\rangle_\mathcal{M},
    \end{equation}

    \begin{lem}
    \label{lem:Eq_u_T}
    For every $(T,x)\in(1,+\infty)\times\Omega$, we have the following:
    \begin{itemize}
        \item for every $(p_T,p_x) \in D^+ u(T,x)$, there exists a curve $\gamma$ realizing $u(T,x)$ in \eqref{eq:u(T,x)_vers2} such that, in the viscosity sense,
        \[
        p_T+H(x,\mathcal M p_x)+\langle \ddot{r}_0(T), \gamma(1)\rangle_\mathcal{M} \le 0;        \]
        \item for every $(p_T,p_x) \in D^- u(T,x)$ and for every $\gamma$ realizing $u(T,x)$ in \eqref{eq:u(T,x)_vers2} it holds, in the viscosity sense,
        \[
        p_T + H(x, \mathcal M p_x) + \langle \ddot{r}_0(T),\gamma(1)\rangle_\mathcal{M} \geq 0.
        \]
    \end{itemize}

    In particular, the function
        $u(T,x)$ defined in \eqref{eq:u(T,x)_vers1} satisfies, in the viscosity sense,
        \begin{equation}
            \label{eq:Eq_u_T}
           -\frac{\partial u}{\partial T}(T,x)=-H(x,\nabla u(T,x)) + \inf_{\gamma \in \mathcal{Y}}\langle \ddot r_0(T),\gamma(T)\rangle_\mathcal{M},
         \end{equation}
        where $\mathcal Y$ is the set of curves $\gamma \in H^1([1,T])$, $\gamma(1)=x$ that realize $u(T,x)$.
    \end{lem}

    \begin{proof}[Proof of Lemma \ref{lem:Eq_u_T}]
To simplify the computations, we consider the formulation of $u(T,x)$ given by the last line in \eqref{eq:u(T,x)_vers2}, so that the end-point of $\gamma$ realizing $u(T,x)$ is $x$.

\smallskip
Take $T>1$, $x\in\Omega$ and $(p_T,p_x)\in D^+ u(T,x)$. Then, for every $z\in\mathcal{X}$, 
\[
\limsup_{h\rightarrow0^+}\frac{u(T-h,x-hz)-u(T,x) + h(p_T + \langle z,p_x\rangle_\mathcal{M})}{h\sqrt{1+\|z\|_\mathcal{M}^2}}\leq 0,
\]
which is equivalent to
\[
\limsup_{h\rightarrow0^+}\frac{u(T-h,x-hz)-u(T,x)}{h}\leq -p_T - \langle z,p_x\rangle_\mathcal{M}.
\]
Setting $\zeta(t) = x + (t-T)z$, with $z\in\mathcal{X}$, we can apply \eqref{eq:dpp} to get
\[
\begin{split}
    u(T,x) &\leq u(T-h,\zeta(T-h)) + \int_{T-h}^{T}L(\zeta(t),\dot{\zeta}(t))\ \ud t - \langle \dot{r}_0(T)-\dot{r}_0(T-h),\gamma_h(1)\rangle_\mathcal{M}\\
    &=u(T-h,x-hz) + \int_{T-h}^{T}L(x+(t-T)z,z)\ \ud t - \langle \dot{r}_0(T)-\dot{r}_0(T-h),\gamma_h(1)\rangle_\mathcal{M},
\end{split}
\]
where $\gamma_h$ realizes the minimum in $u(T-h,\zeta(T-h))$. This implies
\[
\begin{split}
    \limsup_{h\rightarrow 0^+} \frac{u(T-h,x-hz)-u(T,x)}{h} &\geq \lim_{h\rightarrow0^+} -\frac{1}{h} \int_{T-h}^{T}L(x+(t-T)z,z)\ \ud t + \frac{1}{h} \langle \dot{r}_0(T)-\dot{r}_0(T-h),
{\gamma_h}(1)\rangle_\mathcal{M}\\
    &= -L(x,z) + \langle \ddot{r}_0(T),\gamma(1)\rangle_\mathcal{M},
\end{split}
\]
where we used the fact that 
\begin{equation}\label{eq:limit_gamma_h}
    \gamma_h(1)\rightarrow\gamma(1)\text{ as }h\rightarrow0^+,
\end{equation}
with $\gamma$ a minimizer realizing $u(T,x)$. Notice that not every $\gamma$ minimizer is in general obtained as a limit of $\gamma_h$.

\smallskip
We can thus conclude that 
\[
p_T + \langle p_x,z \rangle_\mathcal{M} - L(x,z) + \langle \ddot{r}_0(T),\gamma(1)\rangle_\mathcal{M} \leq 0, \ \mbox{ for all } \ z \in \mathcal X,
\]
which, taking $z=p_x$, yields 
\[
p_T + H(x,\mathcal M p_x) + \langle \ddot{r}_0(T),\gamma(1)\rangle_\mathcal{M} \leq 0,
\]
for every $\gamma \in \mathcal Y$ such that \eqref{eq:limit_gamma_h} holds.


Hence, we have 
\[
p_T + H(x,\mathcal M p_x) + \inf_{\gamma\in\mathcal Y}\langle \ddot{r}_0(T),\gamma(1)\rangle_\mathcal{M} \leq 0.
\]
Now, since $p_x$ is related to $\nabla_{\mathcal M} u = \mathcal M^{-1} \nabla u$, then $u(t,x)$ is a viscosity sub-solution of
\[
\frac{\partial u}{\partial T} +H(x,{\nabla u(T,x)})+\inf_{\gamma\in \mathcal{Y}} \langle \ddot{r}_0(T), \gamma(1)\rangle_{\mathcal M} =0.
\]

To prove that $u$ is a viscosity supersolution, suppose that $\gamma$ is a minimizer for the value function $u(T,x)$ with $x=\gamma(T)$ and $w=\dot{\gamma}(T)$. Then,
\[
\begin{split}
u(T-h,{\gamma}(T-h)) - u(T,x) 
&\leq \int_{1}^{T-h} L(\gamma(t),\dot{\gamma}(t))\ \ud t - \langle \dot{r}_0(T-h), \gamma(1) \rangle_\mathcal{M} - \int_{1}^{T} L(\gamma(t),\dot{\gamma}(t))\ \ud t + \langle \dot{r}_0(T), \gamma(1) \rangle_\mathcal{M}\\
&= -\int_{T-h}^{T} L(\gamma(t),\dot{\gamma}(t))\ \ud t + \langle \dot{r}_0(T) - \dot{r}_0(T-h), \gamma(1) \rangle_\mathcal{M}.
\end{split}
\]
This implies
\[
\begin{split}
    \lim_{h\rightarrow 0^+} \frac{u(T-h,{\gamma}(T-h))-u(T,x)}{h} &\leq \lim_{h\rightarrow0^+} -\frac{1}{h} \int_{T-h}^{T}L(\gamma(t),\dot{\gamma}(t))\ \ud t + \frac{1}{h} \langle \dot{r}_0(T)-\dot{r}_0(T-h),\gamma(1)\rangle_\mathcal{M}\\
    &= -L(x,w) + \langle \ddot{r}_0(T),\gamma(1)\rangle_\mathcal{M}.
\end{split}
\]
Given $(p_T,p_x)\in D^-u(T,x)$, from the fact that $u(T,x)$ is Lipschtiz continuous (see Section \ref{sez_semiconcavity}), it holds
\[
\liminf_{h\rightarrow 0^+} \frac{u(T-h,{\gamma}(T-h))-u(T,x)}{h}\geq -p_T - \langle p_x,w \rangle_\mathcal{M}.
\]
Then
\[
p_T + \langle p_x,w \rangle_\mathcal{M} - L(x,w) + \langle \ddot{r}_0(T),\gamma(1)\rangle_\mathcal{M} \geq 0,
\]
which leads to
\[
p_T + H(x,\mathcal M p_x) + \langle \ddot{r}_0(T),\gamma(1)\rangle_\mathcal{M} \geq 0.
\]
for every $\gamma \in \mathcal Y$

As a consequence, $u(T,x)$ is a viscosity supersolution of 
\[
\frac{\partial u}{\partial T} + H(x,\nabla u(T,x)) + \inf_{\gamma \in \mathcal Y}\langle\ddot{r}_0(T), \gamma(1)\rangle_\mathcal{M} =0.
\]

We have thus proved that the value function $u$ solves, for a fixed $T>1$, the equation
\[
-\frac{\partial u}{\partial T}(T,x) = \frac{1}{2}\|\nabla u(T,x)\|^2_{\mathcal{M}^{-1}} - U(x) + \inf_{\gamma \in \mathcal Y}\langle\ddot{r}_0(T),\gamma(T)\rangle_\mathcal{M}
\]
in the viscosity sense.
\end{proof}

\end{document}
